\section{Open bridges, branch-specific research programme, and promotion criteria}

\subsection{What is now closed at the representation level}

The following chain is mathematically controlled:
\begin{equation}
\boxed{
\vartheta
\xrightarrow{\rm Mellin}
\zeta/\xi
\xrightarrow{\rm Euler\;coordinates}
\{v_p,\log p\}
\xleftarrow{\rm exact\;factorisation}
\text{Collatz prime flow}.
}
\end{equation}
The right-hand arrow is a common-coordinate bridge, not an equality of operators. The exact Collatz branch further provides
\begin{equation}
\boxed{
\mathcal O_m
\longrightarrow
W_s
\longrightarrow
D_C
\longrightarrow
\text{critical conjugate pairs}
\longrightarrow
T^{36}.
}
\end{equation}
The temporal branch provides
\begin{equation}
\boxed{
\Omega
\longrightarrow
\hat H_{ZC}
\longrightarrow
|\Psi(\tau)\rangle
\longrightarrow
A_Z(\tau)
\longrightarrow
\mathcal F_Z(\tau).
}
\end{equation}

\subsection{The remaining hard zeta--Collatz bridge}

The principal mathematical question is now sharply stated:
\begin{quote}
Does a non-arbitrary infinite-dimensional Collatz operator, defined on a mathematically natural function space and using the exact prime/Mellin cocycle or a derived intrinsic weighting, possess a spectral invariant whose analytic continuation and reflection structure are provably related to \(\xi(s)\)?
\end{quote}
A suggestive target has the schematic form
\begin{equation}
\det_{\rm reg}(I-\mathcal L_C(s))
\stackrel{?}{\longleftrightarrow}
\xi(s),
\end{equation}
but the determinant class, function space, regularisation, and weighting must be derived rather than chosen to fit the answer. Until those are fixed, this equation is \status{OPEN}.

\subsection{The remaining time bridge}

The operational phase-dispersion model is closed by definition and standard unitary evolution. The ontological statement
\begin{equation}
\text{``physical time flows because zeta--Schrödinger fuzziness is non-zero''}
\end{equation}
is \status{OPEN}. To promote it beyond interpretation requires a physical model in which:
\begin{enumerate}[leftmargin=2em]
\item the spectral registry is generated by measured or derived physical degrees of freedom rather than selected numerically;
\item the same \(\mathcal F_Z\) predicts independent clock/interference observables;
\item competing standard models are quantitatively compared;
\item the model survives blinded prospective tests.
\end{enumerate}

\subsection{The remaining gravity/Unruh bridge}

The standard relations
\begin{equation}
d\varphi=-E\,d\tau/\hbar,
\qquad
T_U=\hbar a/(2\pi c k_B)
\end{equation}
provide legitimate phase/proper-time and acceleration/frequency structures. A PhaseNav claim that the arithmetic spectrum controls a physical Unruh detector response requires a concrete coupling Hamiltonian and predicted detector transition rate. It is currently \status{OPEN}.

\subsection{The remaining electromagnetic/ZPE bridge}

The classical generator baseline is governed by Maxwell--Faraday dynamics and measured power accounting. The quantum oscillator provides a zero-point scale. The project source equation preserves an intended phase/information modulation. A physical ZPE extraction claim requires, at minimum,
\begin{enumerate}[leftmargin=2em]
\item a Hamiltonian/action that defines the new coupling;
\item a conserved energy accounting including the drive/controller;
\item a predicted signal that differs from conventional induction and losses;
\item calorimetric/electrical/mechanical redundancy;
\item blinded sham/control conditions;
\item independent replication.
\end{enumerate}
No ``watts from \(\Lambda_0\)'' conversion is introduced here.

\subsection{Branch-specific roadmap}

\paragraph{Branch A -- holonomic phase physics.}
Implement a pure connection/holonomy library with dynamic, AB, Berry, and project phase channels typed separately. Test gauge-invariant loop observables. Keep physical units explicit.

\paragraph{Branch B -- theta/zeta analytic controls.}
Maintain high-precision theta/Mellin/\(\xi\) reference identities. Add analytically controlled kernels rather than zero-fitting heuristics.

\paragraph{Branch C -- Collatz prime-flow operator.}
Extend the exact finite cocycle to bounded Koopman matrices on preregistered state spaces. Study stability under boundary changes. Then investigate candidate infinite function spaces and regularised traces/determinants.

\paragraph{Branch D -- temporal synthesis.}
Implement content-addressed spectral registries, direct Fourier synthesis, FFT controls, and NUFFT for non-uniform modes. Benchmark numerical equivalence and precision dependence.

\paragraph{Branch E -- Zeta--Schrödinger fuzziness.}
Implement \(A_Z\), \(\mathcal F_Z\), variance controls, pairwise phase clocks, revivals, and optional open-system extensions. Do not force monotonicity.

\paragraph{Branch F -- \(T^{36}\) runtime binding.}
Map PAIR18 and MODE36 into separate schemas. Bind to the existing v0.27 phase-proper-time state only after semantic equivalence and live \texttt{/dev/shm/ciel\_noema} admission are verified.

\subsection{Promotion ladder}

The recommended evidence ladder is:
\begin{enumerate}[leftmargin=2em]
\item \status{FORMAL\_CLOSED}: identities and definitions proven from stated premises;
\item \status{REFERENCE\_PASS}: independent reference implementation passes deterministic controls;
\item \status{PNCS\_PURE\_PASS}: repository unit tests and provenance close;
\item \status{LIVE\_PHASE\_PASS}: exact \(T^{36}\) execution on the live surface with unchanged semantic projection;
\item \status{SPECTRAL\_THEOREM}: non-fitted operator theorem connecting Collatz and an analytic zeta object;
\item \status{PHYSICAL\_MODEL\_PREDICTION}: quantitative independent prediction beyond the arithmetic/software model;
\item \status{INSTRUMENT\_VALIDATION}: measured prospective signal under controls;
\item \status{INDEPENDENT\_REPLICATION}: external reproduction.
\end{enumerate}
No level inherits the status of the next level automatically.

\subsection{Minimal current synthesis}

The compact v0.2 synthesis is
\begin{align}
\Gamma
&=-\frac{E}{\hbar}d\tau+\frac{q}{\hbar}A_\mu dx^\mu+\mathcal A_a d\lambda^a,\\[1mm]
\vartheta(t)&=t^{-1/2}\vartheta(1/t),\\[1mm]
\xi(s)&=\xi(1-s),\\[1mm]
\zeta(s)&=\prod_p(1-p^{-s})^{-1}\quad(\RePart(s)>1),\\[1mm]
w_s(n)&=\left(\frac{C(n)}n\right)^{-s}(-1)^{n\bmod2},\\[1mm]
W_s(\mathcal O_m)&=\left(\frac{n_m}{n_0}\right)^{-s}(-1)^{N_{\rm odd}},\\[1mm]
W_s&=(-1)^{N_{\rm odd}}\prod_p p^{-s\Delta V_p},\\[1mm]
D_C&=(2\RePart(s)-1)^2\log^2(n_m/n_0),\\[1mm]
|\Psi(\tau)\rangle&=\sum_jc_j e^{-i\omega_j\tau}|j\rangle,\\[1mm]
\mathcal F_Z(\tau)&=1-\left|\sum_j|c_j|^2e^{-i\omega_j\tau}\right|^2.
\end{align}
These formulae define the current research architecture. Only the statements explicitly labelled standard or project-exact are closed.
